\scanpage{186}
On va plutôt essayer de calculer
\[
\tau'-\tau=\operatorname{Tr}xyz-\operatorname{Tr}yxz
=\operatorname{Tr}([x,y]z)=\operatorname{Tr}([y,z]x)=\operatorname{Tr}([z,x].y)
\]
c'est une expression \inserted{multilin.} alternée en $x,y,z$, donc on doit avoir
\[
\tau'-\tau=\deleted{\unsure{\varepsilon}}(a\wedge x\wedge y\wedge z)\otimes e^{-1}
\]
pour $a\in\mathcal{A}$ bien déterminé. Comme cette expression est invariante par $G=\uline{\operatorname{Aut}}\mathcal{A}$, \deleted{on voit que} et que (si $g\in G(-)$)
\[
g(a)\wedge g(x)\wedge g(y)\wedge g(z)=a\wedge x\wedge y\wedge z\quad(=a\wedge g(x)\wedge g(y)\wedge g(z))
\]
on en conclut que $a=ga$ i.e. inv. sous $G$ i.e. $a=\varepsilon.1$. Donc on doit avoir
\[
\tau'-\tau=\varepsilon(1\wedge x\wedge y\wedge z)\otimes e^{-1}
\]
et pour calculer $\varepsilon$, il suffit de prendre $x,y,z$ tels que $1\wedge x\wedge y\wedge z$ inv. i.e. $1,x,y,z$ une base \inserted{de $\mathcal{A}$}. On prendra
\[
x_0=\begin{pmatrix}0&0\\1&0\end{pmatrix},\quad
y_0=\begin{pmatrix}0&1\\0&0\end{pmatrix},\quad
z_0=[x_0,y_0]=\begin{pmatrix}-1&0\\0&1\end{pmatrix}=1-2\alpha
\]
où $\alpha=\begin{pmatrix}1&0\\0&0\end{pmatrix}$
\[
\left(\mathrm{NB}\quad x_0y_0=\begin{pmatrix}0&0\\0&1\end{pmatrix},\quad y_0x_0=\begin{pmatrix}1&0\\0&0\end{pmatrix}\right)
\quad\deleted{\left(x_0y_0z_0=\begin{pmatrix}0&0\\0&1\end{pmatrix}\right)}
\]
On trouve
\[
1\wedge x_0\wedge y_0\wedge z_0=-2\,1\wedge x_0\wedge y_0\wedge\alpha
=-2\alpha'\wedge x_0\wedge y_0\wedge\alpha
=2\deleted{\unsure{x_0\wedge y_0\wedge\alpha\wedge\alpha'}}\inserted{\alpha\wedge\alpha'\wedge x_0\wedge y_0}
\]
où
\[
\alpha=\begin{pmatrix}1&0\\0&0\end{pmatrix},\quad\alpha'=\begin{pmatrix}0&0\\0&1\end{pmatrix}
\]
et \deleted{\unsure{x,y}} $\alpha,\alpha',x_0,y_0$ forment base can. de $M_2(\mathbb{Z})$. On
\margin{185}
